\documentclass[../DoAn.tex]{subfiles}
\begin{document}

Chương này trình bày thiết kế hệ thống EduCenter, bao gồm kiến trúc tổng thể, quy trình nghiệp vụ, thiết kế cơ sở dữ liệu và thiết kế giao diện.

\section{Kiến trúc tổng thể hệ thống}
\label{section:4.1}

Hệ thống EduCenter được xây dựng theo hướng tách lớp rõ ràng giữa giao diện, máy chủ, tầng ca sử dụng và tầng dữ liệu. Về tổng thể, hệ thống gồm bốn khối chính:
\begin{itemize}
    \item Khối giao diện người dùng ở React + TypeScript.
    \item Khối API máy chủ bằng Golang + Gin.
    \item Khối xử lý nghiệp vụ theo hướng Clean Architecture.
    \item Khối lưu trữ dữ liệu PostgreSQL.
\end{itemize}

Từ góc nhìn triển khai, phần giao diện có nhiệm vụ hiển thị màn hình, thu thập dữ liệu đầu vào và gọi API. Phần máy chủ chịu trách nhiệm xác thực, kiểm tra dữ liệu, điều phối ca sử dụng, làm việc với repository và truy xuất dữ liệu từ PostgreSQL. Riêng mô-đun xếp lịch được tổ chức qua lớp trừu tượng \texttt{SchedulingSolver}, cho phép thay thế hoặc so sánh thực nghiệm nhiều thuật toán trên cùng một miền dữ liệu.

\textit{[Chèn Hình 4.1: Kiến trúc tổng thể hệ thống EduCenter]}

\section{Quy trình nghiệp vụ}
\label{section:4.2}

Quy trình nghiệp vụ cốt lõi của hệ thống có thể mô tả theo chuỗi sau:
\begin{enumerate}
    \item Admin tạo các dữ liệu nền như học viên, giáo viên, khóa học, chương trình đào tạo, phòng học và ca học.
    \item Admin mở lớp học, gán chương trình hoặc khóa học, phân công giáo viên và khai báo lịch tuần.
    \item Admin ghi danh học viên vào lớp.
    \item Hệ thống sử dụng thông tin lớp, shift, phòng và giáo viên để tạo phương án xem trước xếp lịch.
    \item Admin xem phương án, đánh giá xung đột và xác nhận.
    \item Hệ thống lưu phương án đã duyệt và sinh buổi học thực tế.
    \item Giáo viên sử dụng buổi học để điểm danh, viết tổng kết và ghi nhận kết quả học tập.
    \item Học viên có thể gửi đơn xin phép nghỉ; quản trị viên hoặc giáo viên xử lý theo quy trình.
\end{enumerate}

\textit{[Chèn Hình 4.2: BPMN hoặc sơ đồ hoạt động cho luồng mở lớp đến xác nhận lưu buổi học]}

\section{Thiết kế cơ sở dữ liệu}
\label{section:4.3}

\subsection{Nguyên tắc đặc tả cơ sở dữ liệu}

Một số quy ước thiết kế dữ liệu đang được dùng thống nhất trong EduCenter:
\begin{itemize}
    \item Khóa chính của hầu hết các bảng sử dụng UUID.
    \item Cơ sở dữ liệu mục tiêu là PostgreSQL.
    \item Các bảng giao dịch quan trọng dùng \texttt{created\_at}, \texttt{updated\_at} để truy vết thay đổi.
    \item Một số bảng dùng \texttt{deleted\_at} để xóa mềm nhằm bảo toàn lịch sử.
\end{itemize}

\subsection{Danh mục các bảng trọng tâm}

\begin{table}[H]
\centering
\caption{Danh mục các bảng trọng tâm của hệ thống}
\label{table:bang4.1}
\begin{tabular}{|l|p{4cm}|p{6cm}|}
\hline
\textbf{Nhóm dữ liệu} & \textbf{Bảng} & \textbf{Vai trò} \\ \hline
Tài khoản và bảo mật & users, user\_otps, password\_resets & Quản lý đăng nhập, kích hoạt tài khoản, quên mật khẩu \\ \hline
Danh mục đào tạo & students, teachers, courses, programs, program\_courses & Quản lý học viên, giáo viên, khóa học và chương trình \\ \hline
Tài nguyên vận hành & rooms, shifts & Chuẩn hóa phòng học và ca học \\ \hline
Tổ chức lớp & classes, class\_schedules, enrollments & Mở lớp, lịch tuần, roster học viên \\ \hline
Vận hành buổi học & lessons, attendances, lesson\_summaries, academic\_records, leave\_requests & Quản lý buổi học và học vụ sau buổi học \\ \hline
\end{tabular}
\end{table}

\textit{[Chèn Hình 4.3: Sơ đồ thực thể - liên kết mức logic của hệ thống EduCenter]}

\subsection{Đặc tả các bảng lõi}

Phần đặc tả dữ liệu được tổng hợp từ ba nguồn chính: các thực thể trong phần máy chủ Golang, các tệp migration PostgreSQL và các ca sử dụng đang vận hành thực tế.

\subsubsection{Bảng users}
\begin{table}[H]
\centering
\caption{Đặc tả bảng users}
\scriptsize
\begin{tabular}{|l|l|l|p{4cm}|}
\hline
\textbf{Trường} & \textbf{Kiểu} & \textbf{Ràng buộc} & \textbf{Ý nghĩa} \\ \hline
id & UUID & PK & Mã định danh tài khoản \\ \hline
code & varchar(50) & unique & Mã nội bộ người dùng \\ \hline
full\_name & varchar(255) & & Họ tên hiển thị \\ \hline
email & varchar(255) & not null, unique & Email đăng nhập chính \\ \hline
password & text & not null & Mật khẩu đã băm \\ \hline
role & varchar(50) & default STUDENT & Vai trò hệ thống \\ \hline
is\_active & boolean & default true & Cờ kích hoạt tài khoản \\ \hline
created\_at, updated\_at & timestamp & audit & Thời gian tạo/cập nhật \\ \hline
deleted\_at & timestamp & soft delete & Phục vụ xóa mềm \\ \hline
\end{tabular}
\end{table}

\subsubsection{Bảng user\_otps}
\begin{table}[H]
\centering
\caption{Đặc tả bảng user\_otps}
\scriptsize
\begin{tabular}{|l|l|l|p{4cm}|}
\hline
\textbf{Trường} & \textbf{Kiểu} & \textbf{Ràng buộc} & \textbf{Ý nghĩa} \\ \hline
id & UUID & PK & Mã OTP record \\ \hline
user\_id & UUID & not null, index & Tham chiếu users.id \\ \hline
otp\_hash & text & not null & Mã OTP đã băm \\ \hline
expired\_at & timestamp & not null & Thời điểm hết hạn \\ \hline
used\_at & timestamp & nullable & Đánh dấu OTP đã dùng \\ \hline
created\_at, deleted\_at & timestamp & audit & Lịch sử xác minh \\ \hline
\end{tabular}
\end{table}

\subsubsection{Bảng password\_resets}
\begin{table}[H]
\centering
\caption{Đặc tả bảng password\_resets}
\scriptsize
\begin{tabular}{|l|l|l|p{4cm}|}
\hline
\textbf{Trường} & \textbf{Kiểu} & \textbf{Ràng buộc} & \textbf{Ý nghĩa} \\ \hline
id & UUID & PK & Mã yêu cầu reset \\ \hline
user\_id & UUID & not null & Tham chiếu users.id \\ \hline
token\_hash & text & not null & Token reset đã băm \\ \hline
expires\_at & timestamp & not null & Hết hạn token \\ \hline
used\_at & timestamp & nullable & Token đã dùng hay chưa \\ \hline
requested\_ip & varchar(45) & & IP gửi yêu cầu \\ \hline
user\_agent & text & & Thông tin thiết bị \\ \hline
created\_at, updated\_at, deleted\_at & timestamp & audit & Truy vết bảo mật \\ \hline
\end{tabular}
\end{table}

\subsubsection{Bảng students}
\begin{table}[H]
\centering
\caption{Đặc tả bảng students}
\scriptsize
\begin{tabular}{|l|l|l|p{4cm}|}
\hline
\textbf{Trường} & \textbf{Kiểu} & \textbf{Ràng buộc} & \textbf{Ý nghĩa} \\ \hline
id & UUID & PK & Mã học viên \\ \hline
code & varchar(50) & unique & Mã học viên \\ \hline
full\_name & varchar(255) & & Họ tên học viên \\ \hline
email & varchar(255) & & Email liên hệ \\ \hline
phone & varchar(20) & & Số điện thoại \\ \hline
guardian\_phone & varchar(20) & & SĐT phụ huynh \\ \hline
grade\_level & varchar(50) & & Khối lớp \\ \hline
status & varchar(50) & default ACTIVE & Trạng thái học viên \\ \hline
date\_of\_birth & timestamp & nullable & Ngày sinh \\ \hline
gender & varchar(20) & & Giới tính \\ \hline
address & text & & Địa chỉ \\ \hline
created\_at, updated\_at, deleted\_at & timestamp & audit / soft delete & Vòng đời hồ sơ \\ \hline
\end{tabular}
\end{table}

\subsubsection{Bảng teachers}
\begin{table}[H]
\centering
\caption{Đặc tả bảng teachers}
\scriptsize
\begin{tabular}{|l|l|p{3cm}|p{3.5cm}|}
\hline
\textbf{Trường} & \textbf{Kiểu} & \textbf{Ràng buộc} & \textbf{Ý nghĩa} \\ \hline
id & UUID & PK & Mã giáo viên \\ \hline
code & varchar(50) & unique & Mã giáo viên \\ \hline
full\_name & varchar(255) & & Họ tên \\ \hline
email, phone & varchar & & Liên hệ \\ \hline
is\_school\_teacher & boolean & default false & GV trường phổ thông hay cộng tác viên \\ \hline
school\_name & varchar(255) & & Trường công tác \\ \hline
employment\_type & varchar(50) & default PART\_TIME & Hình thức làm việc \\ \hline
status & varchar(50) & default ACTIVE & Trạng thái hoạt động \\ \hline
notes & text & & Ghi chú \\ \hline
created\_at, updated\_at, deleted\_at & timestamp & audit / soft delete & Vòng đời hồ sơ \\ \hline
\end{tabular}
\end{table}

\subsubsection{Bảng courses}
\begin{table}[H]
\centering
\caption{Đặc tả bảng courses}
\scriptsize
\begin{tabular}{|l|l|l|p{4cm}|}
\hline
\textbf{Trường} & \textbf{Kiểu} & \textbf{Ràng buộc} & \textbf{Ý nghĩa} \\ \hline
id & UUID & PK & Mã khóa học \\ \hline
code & varchar(50) & unique, not null & Mã khóa học \\ \hline
name & varchar(255) & & Tên khóa học \\ \hline
description & text & & Mô tả nội dung \\ \hline
grade\_level & varchar(50) & & Khối lớp áp dụng \\ \hline
subject & varchar(255) & & Môn học \\ \hline
session\_count & int & & Số buổi học \\ \hline
session\_duration\_minutes & int & & Thời lượng mỗi buổi \\ \hline
total\_hours & numeric(8,2) & & Tổng số giờ học \\ \hline
price & numeric(10,2) & & Học phí \\ \hline
status & varchar(50) & default ACTIVE & Trạng thái \\ \hline
\end{tabular}
\end{table}

\subsubsection{Bảng programs và program\_courses}
\begin{table}[H]
\centering
\caption{Đặc tả bảng programs}
\scriptsize
\begin{tabular}{|l|l|l|p{4cm}|}
\hline
\textbf{Trường} & \textbf{Kiểu} & \textbf{Ràng buộc} & \textbf{Ý nghĩa} \\ \hline
id & UUID & PK & Mã chương trình \\ \hline
code & varchar(50) & unique, not null & Mã chương trình \\ \hline
name & varchar(255) & & Tên chương trình \\ \hline
track & varchar(50) & SUPPORT, BASIC, ADVANCED & Nhánh chương trình \\ \hline
effective\_from, effective\_to & timestamp & nullable & Hiệu lực \\ \hline
created\_by\_id, approved\_by\_id & UUID & nullable & Người tạo/duyệt \\ \hline
published\_at, archived\_at & timestamp & nullable & Xuất bản/lưu trữ \\ \hline
\end{tabular}
\end{table}

Bảng \texttt{program\_courses} là bảng liên kết nhiều-nhiều giữa programs và courses, gồm các trường: id (UUID, PK), program\_id (UUID, not null), course\_id (UUID, not null).

\subsubsection{Bảng rooms}
\begin{table}[H]
\centering
\caption{Đặc tả bảng rooms}
\scriptsize
\begin{tabular}{|l|l|l|p{4cm}|}
\hline
\textbf{Trường} & \textbf{Kiểu} & \textbf{Ràng buộc} & \textbf{Ý nghĩa} \\ \hline
id & UUID & PK & Mã phòng \\ \hline
code & varchar(50) & unique, not null & Mã phòng \\ \hline
name & varchar(255) & not null & Tên phòng \\ \hline
capacity & int & $\geq$ 1 & Sức chứa \\ \hline
address & text & & Địa điểm \\ \hline
\end{tabular}
\end{table}

\subsubsection{Bảng shifts}
\begin{table}[H]
\centering
\caption{Đặc tả bảng shifts}
\scriptsize
\begin{tabular}{|l|l|p{3cm}|p{3.5cm}|}
\hline
\textbf{Trường} & \textbf{Kiểu} & \textbf{Ràng buộc} & \textbf{Ý nghĩa} \\ \hline
id & UUID & PK & Mã ca học \\ \hline
code & varchar(50) & unique, not null & Mã ca \\ \hline
name & varchar(255) & not null & Tên ca \\ \hline
start\_time & varchar(10) & not null & Giờ bắt đầu \\ \hline
end\_time & varchar(10) & not null & Giờ kết thúc \\ \hline
duration\_minutes & int & not null, $\geq$ 1 & Độ dài ca \\ \hline
session\_type & varchar(50) & MORNING, AFTERNOON, EVENING, CUSTOM & Loại ca học \\ \hline
is\_active & boolean & default true & Dùng cho xếp lịch hay không \\ \hline
\end{tabular}
\end{table}

\subsubsection{Bảng classes}
\begin{table}[H]
\centering
\caption{Đặc tả bảng classes}
\scriptsize
\begin{tabular}{|l|l|p{3cm}|p{3.5cm}|}
\hline
\textbf{Trường} & \textbf{Kiểu} & \textbf{Ràng buộc} & \textbf{Ý nghĩa} \\ \hline
id & UUID & PK & Mã lớp học \\ \hline
code & varchar(50) & unique, not null & Mã lớp \\ \hline
name & varchar(255) & not null & Tên lớp \\ \hline
notes & text & & Ghi chú lớp \\ \hline
start\_date & timestamp & not null & Ngày bắt đầu \\ \hline
end\_date & timestamp & nullable & Ngày kết thúc \\ \hline
max\_students & int & & Sĩ số tối đa \\ \hline
status & varchar(50) & default OPEN & Trạng thái lớp \\ \hline
price & numeric(10,2) & & Học phí lớp \\ \hline
program\_id, course\_id, teacher\_id, room\_id & UUID & nullable & Liên kết dữ liệu nền \\ \hline
\end{tabular}
\end{table}

\subsubsection{Bảng class\_schedules}
\begin{table}[H]
\centering
\caption{Đặc tả bảng class\_schedules}
\scriptsize
\begin{tabular}{|l|l|l|p{4cm}|}
\hline
\textbf{Trường} & \textbf{Kiểu} & \textbf{Ràng buộc} & \textbf{Ý nghĩa} \\ \hline
id & UUID & PK & Mã lịch tuần \\ \hline
class\_id & UUID & not null & Tham chiếu classes.id \\ \hline
day\_of\_week & varchar(20) & not null & Thứ học (MONDAY--SUNDAY) \\ \hline
shift\_id & UUID & not null & Tham chiếu shifts.id \\ \hline
room\_id & UUID & nullable & Phòng cố định cho slot \\ \hline
\end{tabular}
\end{table}

\subsubsection{Bảng enrollments}
\begin{table}[H]
\centering
\caption{Đặc tả bảng enrollments}
\scriptsize
\begin{tabular}{|l|l|l|p{4cm}|}
\hline
\textbf{Trường} & \textbf{Kiểu} & \textbf{Ràng buộc} & \textbf{Ý nghĩa} \\ \hline
id & UUID & PK & Mã ghi danh \\ \hline
class\_id & UUID & not null & Tham chiếu lớp \\ \hline
student\_id & UUID & not null & Tham chiếu học viên \\ \hline
status & varchar(50) & default APPLIED & Trạng thái ghi danh \\ \hline
approved\_at & timestamp & nullable & Thời điểm duyệt \\ \hline
rejected\_at & timestamp & nullable & Thời điểm từ chối \\ \hline
\end{tabular}
\end{table}

\subsubsection{Bảng lessons}
\begin{table}[H]
\centering
\caption{Đặc tả bảng lessons}
\scriptsize
\begin{tabular}{|l|l|l|p{4cm}|}
\hline
\textbf{Trường} & \textbf{Kiểu} & \textbf{Ràng buộc} & \textbf{Ý nghĩa} \\ \hline
id & UUID & PK & Mã buổi học \\ \hline
class\_id & UUID & not null & Lớp tương ứng \\ \hline
teacher\_id & UUID & nullable & Giáo viên lesson \\ \hline
date\_start & timestamp & not null & Thời điểm bắt đầu \\ \hline
date\_end & timestamp & not null & Thời điểm kết thúc \\ \hline
room\_id & UUID & nullable & Phòng dạy \\ \hline
notes & text & & Ghi chú \\ \hline
\end{tabular}
\end{table}

\subsubsection{Bảng attendances}
\begin{table}[H]
\centering
\caption{Đặc tả bảng attendances}
\scriptsize
\begin{tabular}{|l|l|p{3cm}|p{3.5cm}|}
\hline
\textbf{Trường} & \textbf{Kiểu} & \textbf{Ràng buộc} & \textbf{Ý nghĩa} \\ \hline
id & UUID & PK & Mã điểm danh \\ \hline
lesson\_id & UUID & not null & Tham chiếu lesson \\ \hline
student\_id & UUID & not null & Tham chiếu học viên \\ \hline
status & int & 1=PRESENT, 2=ABSENT, 3=EXCUSED, 4=LATE, 5=EARLY & Trạng thái chuyên cần \\ \hline
note & text & & Ghi chú \\ \hline
marked\_at & timestamp & nullable & Thời điểm chấm \\ \hline
\end{tabular}
\end{table}

\subsubsection{Bảng lesson\_summaries}
\begin{table}[H]
\centering
\caption{Đặc tả bảng lesson\_summaries}
\scriptsize
\begin{tabular}{|l|l|l|p{4cm}|}
\hline
\textbf{Trường} & \textbf{Kiểu} & \textbf{Ràng buộc} & \textbf{Ý nghĩa} \\ \hline
id & UUID & PK & Mã summary \\ \hline
lesson\_id & UUID & unique, not null & Mỗi lesson tối đa một summary \\ \hline
topic & text & & Chủ đề \\ \hline
lesson\_content & text & & Nội dung đã dạy \\ \hline
class\_feedback & text & & Phản hồi chung \\ \hline
homework & text & & Bài tập \\ \hline
homework\_deadline & timestamp & nullable & Hạn nộp \\ \hline
teacher\_notes & text & & Ghi chú của giáo viên \\ \hline
created\_by\_id & UUID & nullable & Người tạo \\ \hline
\end{tabular}
\end{table}

\subsubsection{Bảng academic\_records}
\begin{table}[H]
\centering
\caption{Đặc tả bảng academic\_records}
\scriptsize
\begin{tabular}{|l|l|p{3cm}|p{3.5cm}|}
\hline
\textbf{Trường} & \textbf{Kiểu} & \textbf{Ràng buộc} & \textbf{Ý nghĩa} \\ \hline
id & UUID & PK & Mã record \\ \hline
lesson\_summary\_id & UUID & not null & Summary cha \\ \hline
student\_id & UUID & not null & Học viên được đánh giá \\ \hline
homework\_completed & boolean & default false & Đã hoàn thành bài tập \\ \hline
homework\_score & numeric(5,2) & & Điểm bài tập \\ \hline
attitude\_rating & int & & Đánh giá thái độ \\ \hline
participation\_score & numeric(5,2) & & Điểm tham gia \\ \hline
personal\_comment & text & & Nhận xét cá nhân \\ \hline
total\_score & numeric(5,2) & & Tổng điểm \\ \hline
is\_completed & boolean & default false & Đã chốt record \\ \hline
\end{tabular}
\end{table}

\subsubsection{Bảng leave\_requests}
\begin{table}[H]
\centering
\caption{Đặc tả bảng leave\_requests}
\scriptsize
\begin{tabular}{|l|l|p{3cm}|p{3.5cm}|}
\hline
\textbf{Trường} & \textbf{Kiểu} & \textbf{Ràng buộc} & \textbf{Ý nghĩa} \\ \hline
id & UUID & PK & Mã đơn \\ \hline
student\_id & UUID & not null & Học viên gửi đơn \\ \hline
leave\_type & varchar(50) & LEAVE, LATE, EARLY & Loại đơn \\ \hline
apply\_date & timestamp & not null & Ngày áp dụng \\ \hline
late\_minutes & int & dùng cho LATE & Số phút đi muộn \\ \hline
early\_minutes & int & dùng cho EARLY & Số phút về sớm \\ \hline
reason & text & not null & Lý do \\ \hline
documents & text[] & nullable & Minh chứng đính kèm \\ \hline
class\_id, lesson\_id & UUID & nullable & Lớp/lesson liên quan \\ \hline
status & varchar(50) & default PENDING & Trạng thái xử lý \\ \hline
approved\_by\_id & UUID & nullable & Người duyệt \\ \hline
rejection\_reason & text & & Lý do từ chối \\ \hline
\end{tabular}
\end{table}

\subsection{Đặc tả các bảng hỗ trợ và mở rộng}

Nhóm bảng hỗ trợ bao gồm: \texttt{consultations} (yêu cầu tư vấn đầu vào), \texttt{objectives} (mục tiêu đào tạo) và \texttt{outcomes} (chuẩn đầu ra). Ba bảng này bổ sung cho hệ thống ở hai chiều: giai đoạn trước vận hành lớp học và chiều sâu dữ liệu đào tạo.

\subsection{Giá trị mặc định và domain dữ liệu}

\begin{table}[H]
\centering
\caption{Giá trị mặc định và domain dữ liệu quan trọng}
\label{table:bang4.2}
\scriptsize
\begin{tabular}{|l|p{5cm}|p{4cm}|}
\hline
\textbf{Trường} & \textbf{Giá trị mặc định / domain} & \textbf{Ghi chú} \\ \hline
users.role & Mặc định STUDENT & Vai trò khi đăng ký mới \\ \hline
students.status & Mặc định ACTIVE & Học viên hoạt động \\ \hline
teachers.employment\_type & PART\_TIME, FULL\_TIME; mặc định PART\_TIME & Loại hình cộng tác \\ \hline
shifts.session\_type & MORNING, AFTERNOON, EVENING, CUSTOM & Loại ca học \\ \hline
classes.status & OPEN, CLOSED, CANCELLED; mặc định OPEN & Xếp lịch chỉ dùng OPEN \\ \hline
leave\_requests.status & PENDING, APPROVED, REJECTED & Trạng thái xử lý \\ \hline
attendances.status & 1=PRESENT, 2=ABSENT, 3=EXCUSED, 4=LATE, 5=EARLY & Domain nội bộ \\ \hline
\end{tabular}
\end{table}

\section{Thiết kế giao diện}
\label{section:4.4}

Phần giao diện của hệ thống được xây dựng theo hướng phục vụ thao tác quản trị rõ ràng, bám theo từng nhóm chức năng.

\subsection{Giao diện xác thực và tài khoản}

\textit{[Chèn Hình 4.4: Giao diện đăng nhập]}

\subsection{Giao diện quản lý học viên và giáo viên}

\textit{[Chèn Hình 4.5: Giao diện quản lý học viên]}

\textit{[Chèn Hình 4.6: Giao diện quản lý giáo viên]}

\subsection{Giao diện quản lý khóa học, chương trình và lớp học}

\textit{[Chèn Hình 4.7: Giao diện quản lý khóa học]}

\textit{[Chèn Hình 4.8: Giao diện quản lý chương trình đào tạo]}

\textit{[Chèn Hình 4.9: Giao diện quản lý lớp học]}

\subsection{Giao diện quản lý phòng học và ca học}

\textit{[Chèn Hình 4.10: Giao diện quản lý phòng học]}

\textit{[Chèn Hình 4.11: Giao diện quản lý ca học]}

\subsection{Giao diện xếp lịch và buổi học}

\textit{[Chèn Hình 4.12: Giao diện bước 1 chọn lớp hoặc giáo viên để xếp lịch]}

\textit{[Chèn Hình 4.13: Giao diện xem trước lịch học trên lưới thời gian]}

\textit{[Chèn Hình 4.14: Giao diện danh sách buổi chưa xếp được và xung đột]}

\textit{[Chèn Hình 4.15: Giao diện xác nhận lưu để sinh các buổi học thực tế]}

\end{document}
